\documentclass[11pt]{article}
\usepackage[margin=0.4in]{geometry}
\usepackage{hyperref}
\usepackage{enumitem}
\usepackage{titlesec}
\usepackage{parskip}
\usepackage{booktabs}

\titlespacing*{\section}{0pt}{8pt}{4pt}
\titlespacing*{\subsection}{0pt}{6pt}{3pt}

\setlength{\parskip}{4pt}
\setlength{\parindent}{0pt}

\hypersetup{colorlinks=true, linkcolor=blue, urlcolor=blue, citecolor=blue}

\begin{document}

\begin{center}
    {\large \textbf{Project Proposal -- Empirical Analysis of Saddle-Point Escaping Optimizers on Non-Convex ML Loss Landscapes}}\\[4pt]
    Shyam Patadia \quad \texttt{spatadia@wpi.edu}\\
    MA 551 Computational Statistics
\end{center}

\section{Introduction}

Training a machine learning model is fundamentally an optimization problem: minimize a loss function $f(\theta)$ over a high-dimensional, non-convex parameter space. In practice, neural network loss landscapes are riddled with saddle points -- critical points where the gradient vanishes but the point is neither a minimum nor a maximum -- and flat plateau regions where gradients are near-zero and learning stalls. In high dimensions, saddle points vastly outnumber local minima, making them the dominant obstacle in deep learning optimization.

A recent algorithm, Steepest Perturbed Gradient Descent (SPGD) \cite{vahedi2024spgd}, proposes a principled escape mechanism: every $\textit{IterP}$ iterations, generate $N_P$ uniformly perturbed candidate solutions around the current point and select the one yielding the steepest loss decrease. This combines the directed efficiency of gradient descent with a Monte Carlo exploration strategy. However, SPGD has only been evaluated on mathematical benchmark functions and a 3D engineering problem; its behavior on ML loss landscapes is entirely unexplored. This project provides the first systematic empirical study of SPGD on ML tasks, comparing it against established optimizers under controlled conditions.

\section{Problem Description}

Standard optimizers -- SGD, Adam, and their variants -- rely solely on gradient information. Near saddle points, gradients approach zero, causing these methods to stall for many iterations with negligible progress. Perturbed Gradient Descent (PGD) \cite{jin2017escape} adds random noise to escape such regions but treats all perturbation directions equally. SPGD's steepest-selection rule is more principled: it evaluates $N_P$ candidates and selects the most promising one. Whether this selection strategy provides a meaningful advantage over simpler perturbation approaches -- and under what conditions -- is an open empirical question with direct implications for ML practitioners choosing optimizers for non-convex problems.

\section{Objectives}

\begin{enumerate}[leftmargin=*, itemsep=2pt, topsep=2pt]
    \item Implement SPGD as a PyTorch \texttt{Optimizer} subclass and validate it on non-convex benchmark functions (Rastrigin, Ackley, Rosenbrock) at $d \in \{2, 10, 50\}$.
    \item Compare SPGD against SGD, Adam, PGD, and a random-perturbation baseline across four datasets spanning synthetic, tabular, and vision settings.
    \item Diagnose saddle-point behavior through gradient-norm stagnation episodes and escape-time measurements to determine when and why SPGD outperforms competing methods.
    \item Characterize SPGD's sensitivity to $N_P$ and $\textit{IterP}$ via ablation on the Two Moons dataset.
\end{enumerate}

\section{Datasets and Experiments}

\begin{table}[h]
\centering
\small
\begin{tabular}{@{}llllr@{}}
\toprule
\textbf{Dataset} & \textbf{Type} & \textbf{Size} & \textbf{What it isolates} & \textbf{Est.\ compute} \\
\midrule
Rastrigin / Ackley / Rosenbrock & Synthetic & Generated & Benchmark validation, visual trajectories & $<$1 min \\
Two Moons (\texttt{sklearn}) & Synthetic 2D & 1{,}000 & Visual saddle escape, loss landscape plot & $<$5 min \\
OpenML-CC18 tabular subset & Real tabular & 500--50k & Real-world generalization, 5 datasets & $<$20 min \\
CIFAR-10 (\texttt{torchvision}) & Real images & 60{,}000 & Scale stress test, practical convergence & $\sim$2 hrs \\
\bottomrule
\end{tabular}
\end{table}

\subsection*{Experiment 1 -- Benchmark Function Validation (Week 1)}

SPGD is validated on Rastrigin ($f(x) = An + \sum[x_i^2 - A\cos(2\pi x_i)]$), Ackley, and Rosenbrock functions at $d \in \{2, 10, 50\}$ with 30 random initializations per optimizer. These functions have known global optima and well-characterized saddle/local-minima structure. Metrics: final function value, iterations to convergence, and success rate (fraction of runs reaching within $\epsilon = 0.01$ of the global optimum). 2D trajectory plots visualize SPGD escaping local minima that trap vanilla GD.

\subsection*{Experiment 2 -- Two Moons: Visual Saddle Analysis (Week 2)}

A small MLP (2$\to$32$\to$32$\to$1, no batch normalization) is trained on Two Moons for 2{,}000 steps across all five optimizers with 5 random seeds. The loss landscape is visualized by sweeping two weight dimensions and plotting the surface alongside optimizer trajectories. Stagnation episodes -- consecutive steps where $\|\nabla f\| < 10^{-4}$ -- are counted per optimizer. This is the primary visual experiment, producing the most interpretable figure in the paper.

\subsection*{Experiment 3 -- OpenML Tabular: Real-Data Generalization (Week 2)}

Five tabular classification datasets are drawn from the OpenML-CC18 benchmark suite (\texttt{pip install openml}), selected with $n \in [500, 50{,}000]$ rows and $p \in [10, 100]$ features. A 3-layer MLP is trained on each dataset for all five optimizers across 3 seeds. Metrics: training loss convergence, stagnation episode count, and test accuracy. This experiment tests whether SPGD's advantage observed on synthetic problems transfers to real heterogeneous tabular data.

\subsection*{Experiment 4 -- CIFAR-10: Scale Stress Test (Weeks 3--4)}

A ResNet-18 (11M parameters, no batch normalization in early layers to preserve saddle structure) is trained on CIFAR-10 for 50 epochs across all five optimizers with 3 seeds. This is the anchor experiment. Metrics: training loss curve, gradient norm over time, stagnation episode count, final test accuracy, and sharpness at convergence (loss increase under small weight perturbation $\delta = 0.01$, following the flat-minima hypothesis). Total GPU time: approximately 2 hours on an RTX 4050.

\subsection*{Experiment 5 -- Ablation on $N_P$ and $\textit{IterP}$ (Week 4)}

A $4 \times 4$ grid over $N_P \in \{1, 5, 10, 20\}$ and $\textit{IterP} \in \{5, 10, 50, 100\}$ is run on Two Moons (16 combinations $\times$ 5 seeds = 80 runs, each under 5 seconds). The result is a heatmap of final test accuracy and stagnation count, identifying the practical operating regime for SPGD.

\section{Approach and Methodology}

\textbf{Implementation.} SPGD is implemented as a PyTorch \texttt{Optimizer} subclass. Every $\textit{IterP}$ steps, $N_P$ candidates are sampled by adding uniform perturbations $\delta \sim \text{Uniform}(-r, r)^d$ to current parameters. The candidate with the largest loss decrease is accepted; otherwise the standard gradient step is taken. The random-perturbation baseline uses the same mechanism but selects the candidate randomly rather than by steepest descent, isolating the contribution of SPGD's selection rule.

\textbf{Diagnostics.} Stagnation episodes (steps with $\|\nabla f\| < \varepsilon$) and escape times are logged at zero cost inside the training loop. These diagnostics directly measure saddle-point behavior rather than inferring it from loss curves alone.

\textbf{Reproducibility.} All experiments use fixed random seeds. Total compute across all experiments is under 3 hours on an RTX 4050.


\begin{thebibliography}{9}
\bibitem{vahedi2024spgd}
A.~M.~Vahedi and H.~T.~Ilies, ``SPGD: Steepest Perturbed Gradient Descent Optimization,'' \textit{arXiv:2411.04946}, 2024.

\bibitem{jin2017escape}
C.~Jin, R.~Ge, P.~Netrapalli, S.~M.~Kakade, and M.~I.~Jordan, ``How to Escape Saddle Points Efficiently,'' \textit{ICML}, 2017.
\end{thebibliography}

\end{document}